% BHU PYQ -- Transcribed & Corrected
% Paper: CHEMJ21 / CHEMN21: Basic Concepts of Chemistry-II
% Exam: B.Sc. (Hons) Semester III (NEP) Examination 2024-25
% Department: Chemistry

\documentclass[11pt,a4paper]{article}
\usepackage[a4paper,top=2.5cm,bottom=2.5cm,left=2.8cm,right=2.8cm,headheight=14pt]{geometry}
\usepackage{amsmath,amssymb,chemformula}
\usepackage[T1]{fontenc}
\usepackage[utf8]{inputenc}
\usepackage{lmodern,microtype,fancyhdr,enumitem,hyperref,lastpage,parskip,graphicx}

\hypersetup{
    colorlinks=true,
    urlcolor=black,
    linkcolor=black,
    pdftitle={CHEMJ21 / CHEMN21: Basic Concepts of Chemistry-II -- BHU Sem III 2024-25}
}

\pagestyle{fancy}
\fancyhf{}
\renewcommand{\headrulewidth}{0.4pt}
\renewcommand{\footrulewidth}{0.4pt}
\fancyhead[L]{\small\textit{Banaras Hindu University}}
\fancyhead[R]{\small\textit{CHEMJ21 / CHEMN21: Basic Concepts of Chemistry-II}}
\fancyfoot[C]{\small Page \thepage\ of \pageref{LastPage}}

\newlist{parts}{enumerate}{2}
\setlist[parts,1]{label=(\alph*),leftmargin=*,itemsep=4pt,topsep=2pt}
\setlist[parts,2]{label=(\roman*),leftmargin=*,itemsep=2pt,topsep=2pt}
\newcommand{\pts}[1]{\hfill\textit{[#1]}}

\begin{document}

\begin{center}
  {\large\bfseries BANARAS HINDU UNIVERSITY}\\[2pt]
  {\normalsize B.Sc. (Hons) Semester III (NEP) Examination 2024-25}\\[6pt]
  \rule{0.8\linewidth}{0.5pt}\\[6pt]
  {\large\bfseries CHEMISTRY}\\[4pt]
  {\large\bfseries Paper No. CHEMJ21 / CHEMN21 : Basic Concepts of Chemistry-II}\\[4pt]
  {\normalsize Time: 3 Hours\hfill Maximum Marks: 70}
\end{center}

\rule{\linewidth}{0.4pt}

\smallskip
\noindent(Write your Roll No. at the top immediately on the receipt of this question paper)

\smallskip
\noindent\textit{Attempt five questions including Question 1 which is compulsory. Attempt at least one question from Questions 2 and 3. The figures in the right-hand margin indicate marks.}

\bigskip

\begin{enumerate}[label=\textbf{\arabic*.}]

    \item Attempt all of the following:
    \begin{parts}
        \item ``Carbon possesses only two unpaired spins in its ground state yet it forms four bonds.'' Why?\pts{2}
        \item Comment upon stability of $\ch{Be2}$ as per molecular orbital theory.\pts{2}
        \item ``The H--C--F bond angle is slightly smaller than H--C--H in \ch{CH3F}.'' Explain this observation in terms of Bent's rule.\pts{1}
        \item Explain why the phenols are stronger acids than alcohols but weaker than carboxylic acids.\pts{2}
        \item Why does the fully-eclipsed Newman projection of n-butane have maximum torsional strain?\pts{2}
        \item Why does only the $\text{E}2$ elimination show the kinetic isotope effect while $\text{E}1$ does not?\pts{2}
        \item In Williamson's ether synthesis, why is a primary alkyl halide used instead of a secondary or tertiary one?\pts{2}
        \item How many transition states exist in the $\text{S}_\text{N}1$ reaction?\pts{1}
    \end{parts}

    \bigskip

    \item Answer the following:
    \begin{parts}
        \item Explain in brief the Valence Bond Theory of metallic bonding through a suitable example.\pts{4}
        \item ``The boiling point of HF is much higher than that of HCl, HBr, and HI, while in aqueous state it is the poorest acid among all.'' Why?\pts{4}
        \item Discuss the structure and bonding of either Methyl Lithium or Diborane.\pts{6}
        \item Draw the crystal structure of Rock salt ($\ch{NaCl}$) and mention the coordination numbers of sodium and chloride ions.\pts{4}
        \item Why is the Si--F bond stronger than the corresponding C--F bond?\pts{2}
    \end{parts}

    \bigskip

    \item Answer the following:
    \begin{parts}
        \item Let A and B be two atoms with wave functions for 1s orbitals as $\psi_A$ and $\psi_B$, respectively. Write the molecular orbital wave functions using these atomic orbital wave functions $\psi_A$ and $\psi_B$. Write expressions using molecular orbitals for the probability of finding the electron density in a unit volume and interpret the meaning of each term in the expressions.\pts{6}
        \item ``The compounds trimethylamine $\ch{(CH3)3N}$ and trimethylsilylamine $\ch{(Me3Si)3N}$ have similar formulae but have pyramidal and trigonal planar structures, respectively.'' Why?\pts{4}
        \item Explain in brief the Molecular Orbital Theory of metallic bonding through a suitable example.\pts{4}
        \item Comment upon the variation of bond angles in $\ch{CH4}$, $\ch{NH3}$, and $\ch{H2O}$ in the light of VSEPR theory.\pts{4}
        \item Explain the defects noted in extending homonuclear diatomic molecular orbital concepts to the heteronuclear diatomic molecule CO using an energy level diagram.\pts{2}
    \end{parts}

    \bigskip

    \item Answer the following:
    \begin{parts}
        \item Write down the product(s) of the following reactions:\pts{8}
        \begin{parts}
            \item $\ch{H3C-C \equiv CH} \xrightarrow[\text{dil. } \ch{H2SO4}]{\ch{HgSO4}} \textbf{?}$
            \item Cyclohexene $\xrightarrow[\text{2. } \ch{H^+}]{\text{1. MCPBA}} \textbf{?}$
            \item $n\,\ch{H3C-CH=CH2} \xrightarrow{\text{Benzoyl peroxide}} \textbf{?}$
        \end{parts}
        \item Arrange the following in increasing order of basicity with a suitable explanation:\pts{4}
        \begin{parts}
            \item $\ch{Ph-NH2}$
            \item $\ch{Ph2NH}$
            \item $\ch{Ph3N}$
        \end{parts}
        \item Predict the E and Z nomenclature of the following compounds:\pts{2}
        \begin{parts}
            \item 1-bromo-1-chloroprop-1-ene
            \item 1-chloro-2-deuterioethene derivatives
        \end{parts}
    \end{parts}

    \bigskip

    \item Answer the following:
    \begin{parts}
        \item Write down the major products formed in the following reaction with a suitable mechanistic explanation:\pts{8}\\[4pt]
        $\ch{(CH3)2CH-CH(Br)-CH3} \xrightarrow{\ch{CH3OK} / \ch{MeOH}} \textbf{?}$
        \item What is the intermediate involved in the $\text{E}1\text{cB}$ mechanism?\pts{2}
        \item Identify the products formed in the following reaction with a suitable mechanism:\pts{4}\\[4pt]
        $\ch{CH3-O-C(CH3)3 + HI} \longrightarrow \textbf{?}$
    \end{parts}

    \bigskip

    \item Answer the following:
    \begin{parts}
        \item Predict the R and S nomenclature of the following chiral molecules:\pts{3}
        \begin{parts}
            \item Glyceraldehyde derivative
            \item L-amino acid derivative
        \end{parts}
        \item How do you prepare the following from phenylmagnesium iodide ($\ch{PhMgI}$)?\pts{6}
        \begin{parts}
            \item Benzenesulfonic acid ($\ch{PhSO3H}$)
            \item Phenol ($\ch{PhOH}$)
            \item Benzyl alcohol ($\ch{PhCH2OH}$)
        \end{parts}
        \item How do you prepare the following from ethyl acetoacetate?\pts{5}
        \begin{parts}
            \item A 1,4-diketone
            \item A 1,6-diketone
        \end{parts}
    \end{parts}

    \bigskip

    \item Answer the following:
    \begin{parts}
        \item Identify the products formed in the following reaction with a suitable mechanism:\pts{4}\\[4pt]
        $\ch{CH2=CH-CH=CH2 + HBr} \longrightarrow \textbf{?}$
        \item How do you prepare the following from diethyl malonate?\pts{4}
        \begin{parts}
            \item Glutaric acid
            \item Dimethyl acetic acid
        \end{parts}
        \item Discuss the $\text{S}_\text{N}1$ reaction mechanism with a suitable example.\pts{4}
        \item Distinguish between primary and secondary alcohols by the substitution method.\pts{2}
    \end{parts}

\end{enumerate}

\end{document}
